%% kapa: exam=1 round=34 year=2023 period=1 form=A subject=civil_law subject_ko=민법 seq=36 number=36
%% source: Q-Net 감정평가사 기출문제 (cst00309) artlSeq=5211789 file=raw/1008/5211789/2023년도 제34회 감정평가사 1차시험 1교시.pdf
%% conversion: kordoc 4.24.51 -> tools/gen1.py (auto); review=proofread-sonnet; flags=none
\begin{question}{36}
유치권에 관한 설명으로 옳지 않은 것은? (다툼이 있으면 판례에 따름)
\begin{choices}
  \item 유치물의 소유자가 변동된 후 유치권자가 유치물에 관하여 새로이 유익비를 지급하여 가격증가가 현존하는 경우, 유치권자는 그 유익비를 피보전채권으로 하여서도 유치권을 행사할 수 있다.
  \item 다세대주택의 창호공사를 완성한 하수급인이 공사대금채권 잔액을 변제받기 위하여 그 중 한 세대를 점유하는 유치권 행사는 인정되지 않는다.
  \item 수급인의 재료와 노력으로 건물을 신축한 경우, 특별한 사정이 없는 한 그 건물에 대한 수급인의 유치권은 인정되지 않는다.
  \item 유치권의 목적이 될 수 있는 것은 동산, 부동산 그리고 유가증권이다.
  \item 유치권자가 유치물에 대한 보존행위로서 목적물을 사용하는 것은 적법하다.
\end{choices}
\end{question}
\answer{2}
